\section{Why the Solution Should Work}
\label{sec:justification}

% Q4: Why do you think your solution will work?

\subsection{Theoretical Motivation}

Our approach is grounded in several key insights:

\paragraph{Topology is Learnable}
Critical points in scalar fields have distinctive local signatures---the gradient vanishes and the Hessian matrix determines the type. An RL agent with access to local gradient and curvature information can learn to identify and prioritize these regions.

\paragraph{Diffusion Models Learn Physical Priors}
Diffusion models have shown remarkable ability to learn complex data distributions. When trained on scientific data, they implicitly learn physical constraints and correlations. This enables high-quality reconstruction even from heavily compressed representations.

\paragraph{RL Can Learn Non-Obvious Strategies}
Unlike hand-crafted heuristics, reinforcement learning can discover compression strategies that are not immediately intuitive to human designers. The agent may learn to:
\begin{itemize}
    \item Preserve certain saddle points that are critical for filament structure
    \item Aggressively compress flat regions while maintaining precision near extrema
    \item Adapt strategy based on local data characteristics
\end{itemize}

\subsection{Connection to Prior Work}

Our approach builds on established foundations:

\begin{itemize}
    \item \textbf{Adaptive compression}: Prior work shows that spatially-varying compression can outperform uniform approaches
    \item \textbf{Conditional generation}: Diffusion models conditioned on partial information (inpainting, super-resolution) achieve high quality
    \item \textbf{RL for optimization}: RL has been successfully applied to learn policies for complex optimization problems
\end{itemize}

\subsection{Expected Outcomes}

Based on our analysis, we expect:

\begin{enumerate}
    \item \textbf{Higher critical point preservation} at equivalent compression ratios compared to uniform methods
    \item \textbf{Interpretable policies} where the agent allocates more bits to topologically important regions
    \item \textbf{Graceful degradation} where topology preservation remains high even at aggressive compression levels
\end{enumerate}

\subsection{Potential Challenges}

We acknowledge several challenges:

\begin{itemize}
    \item Training stability of RL with sparse topology rewards
    \item Computational cost of diffusion model inference
    \item Generalization across different scientific domains
\end{itemize}

These challenges inform our experimental design and evaluation criteria.
